\section{気候評価と排出管理は対象量を分ける}
\label{theory:env:climateaction}
\label{theory:env:ar6}
\label{theory:env:blue}
オゾン層破壊を抑える代替物質でも、強い温室効果を持つ場合がある。オゾン層破壊係数と地球温暖化係数は異なる指標である。RE100は電力調達、TCFDは財務情報の開示、ZEH・ZEBは建築の省エネルギー性能と収支を扱う。制度名の共通する「気候」という語だけで対象量を同一視しない。

AR6 WGIが示す2011〜2020年と1850〜1900年の世界平均地表温度差は約1.1℃である。将来の排出シナリオや陸域の値をこの期間平均へ代入しない。影響と脆弱性、排出削減はそれぞれ別の作業部会で評価される。\cite{env60_ar6}

日本は2024年提出インベントリで海草・海藻藻場の吸収量を合わせて国連へ報告した。対象年度は2022年度であり、報告年と同じではない。海洋全体の炭素循環を全て確定したという意味にはしない。\cite{env60_blue}

\section{再エネの比率と風のエネルギー流束}
\label{theory:env:renewables}
\label{theory:env:energy2020}
\label{theory:env:oil2019}
発電量の再エネ比と一次エネルギー供給の再エネ比は分母が違う。水力を含むか、未活用エネルギーを加えるか、対象が暦年か年度かも確かめる。2023年度の発電量における水力を含む再エネ比は22.9\%であり、2018年度の水力を除く比へ転用しない。\cite{env60_re2023}

設備費の単純平均は、容量で重み付けした平均や価格算定の想定値とは異なる。万円/kWは設備容量当たりの費用、円/kWhは発電量当たりの費用である。住宅用と事業用の費用も分ける。個別年の上下を確認せずに「全規模で毎年低下」と一般化しない。\cite{env60_pv2024}

面積 \(A\) を通過する風の質量流量は \(\rho Av\)、単位質量の運動エネルギーは \(v^2/2\) だから、
\[
 P_{\rm wind}=\frac12\rho Av^3.
\]
これは風が運ぶ流束であり、実際の機械出力には空力係数と効率が掛かる。一定質量の運動エネルギーが風速二乗で増えることと混同しない。

国内導入容量が増えても国内製造品の出荷比率が増えるとは限らない。原油の中東依存度は数量の地域比、輸入額比は金額の品目比である。石油とLNGの調達構成、資源量と導入容量も別に比較する。

\section{水素は物性・反応熱・全過程を区別する}
\label{theory:env:hydrogen}
水素の常圧での液化温度は約−253℃、液体体積は常温気体の概数1/800である。水素は単位質量当たりの発熱量が大きいが、体積当たりでも常にガソリンより大きいとはいえない。

水素による鉄鉱石還元は加熱を要する吸熱側の技術として扱われる。\cite{env60_hydrogen}
厳密な反応熱の計算には酸化物、反応温度、生成水の相を指定する。反応でCO\(_2\)を直接発生しなくても、水素製造と電力供給を含む全過程の排出がゼロとは限らない。一次資源から水素へ変換するため、水素は二次エネルギーとして考える。

\section{湿地と里山の保全は同じ管理を要求しない}
\label{theory:env:biodiversity2023}
\label{theory:env:treaties}
ラムサール条約は湿地の保全と賢明な利用を扱う。ワシントン条約は種の国際取引、モントリオール議定書は対象物質の生産・消費、バーゼル条約は有害廃棄物等の越境移動、名古屋議定書は遺伝資源の取得と利益配分を扱う。対象と行為を一組で覚える。\cite{env60_ramsar}

生物多様性の損失には開発・改変だけでなく、人の利用・管理の低下も関わる。里地里山の二次的自然は一定の人間活動で維持されるため、利用停止が生息環境を変えることがある。開発を止めれば全ての生態系が自動的に回復するという単一の説明にしない。\cite{env60_biodiversity}

\section{世帯排出の図は燃料量の図ではない}
\label{theory:env:house2022}
2022年度の全国世帯平均の指定図では、電気1.74、都市ガス0.38、LPガス0.15、灯油0.32、合計2.59（t-CO\(_2\)/世帯・年）となる。\cite{env60_house2022}
これらは同じ体積又は熱量の燃料使用量ではない。電気の排出係数には発電に伴う排出の帰属が含まれる。年ごとの変化は使用量、気温、機器、係数等を分けて考え、積み上げ棒の大小だけから一要因を断定しない。

\section{政策の基本的視点と不確実性への対応}
\label{theory:env:energyplan7}
エネルギー政策は安全性を大前提に、安定供給、経済効率性、環境適合性を考えるS+3Eを基本にする。将来需要や技術の進歩には不確実性があり、使える技術の選択肢を幅広く追求する。多様な選択肢の評価と、最初から限定する方針を同義にしない。\cite{env60_plan7}
削減目標では基準年度と目標年度を併記し、60\%や73\%という数字だけを覚えない。

\section{基本計画は期と特徴句を対応させる}
\label{theory:env:scienceplans}
第1期1996のポストドクター等1万人支援、第2期2001の社会のため・社会の中の科学技術、第3期2006の倫理的・法的・社会的課題、第4期2011の震災復興、第5期2016のSociety 5.0、第6期2021の総合知という対応を整理する。\cite{env60_plan2,env60_plan3}
特徴句の対応は、その概念がそれ以前には存在しなかったという主張ではない。法律の改正年、名称変更の施行年、五年計画の開始年を区別する。

\section{制定・発明・受賞の年を混ぜない}
\label{theory:env:safetyhistory}
\label{theory:env:nobeljapan}
工場法制定1911、全国安全週間初回1928、工業標準化法制定1949、労働安全衛生法制定1972、製造物責任法制定1994を順序として整理する。工場法施行1916、製造物責任法施行1995は別の出来事である。\cite{env60_safety}

1949年の物理学賞は湯川秀樹の中間子予言に対する賞であり、朝永振一郎の量子電気力学に対する1965年の賞とは異なる。\cite{env60_yukawa}
特許第1号の堀田の錆止塗料は1885年7月1日に出願されたものである。出願日を登録日へ置き換えない。\cite{env60_patent}
発明、実験による確認、量産、販売、受賞を同じ「成立年」で一覧にせず、問題が指定する業績の段階を確かめる。
